\subsection{Energy on a metric graph}
\label{sec:ch26-ex-graph}

We retain the original BFJ hypotheses of \eqref{eq:ch20-standing},
including the normalization \(-\log|\pi|=1\).  In particular, the
reference \(\omega\) may be an arbitrary real semipositive closed model
form with ample generic class; it need not be the Chern class of one
metrized line bundle.  The curve need only be connected, not geometrically
connected.  For an edge \(e\) of a regular semistable model, write
\[
 c_e=[\kappa(z_e):k],\qquad k=\widetilde K,
\]
where \(z_e\) is its node.  Slopes are taken in the analytic modulus
coordinate.  The same weight \(c_e\) is used at both ends of the edge
and on all its subdivisions; it is not the degree of the constant field
of either adjacent component.

\begin{lemma}
\label{lem:ch26-weighted-graph-real-bfj}
Let \(\mathcal X\) be a regular semistable model of the smooth connected
projective curve \(X\), determining \(\omega\), and let
\(r:X^{\An}\to\Gamma\) be its skeletal retraction.  Put
\(\mu_0=\MA_\omega(0)\), of mass \(V=(\{\omega\})>0\).
Use the residue-weighted outgoing-slope Laplacian
\(\Delta_\Gamma\) of \eqref{eq:ch31-graph-laplacian}.
For every positive Radon measure \(\mu\) on \(\Gamma\) of mass \(V\),
the normalized continuous graph solution
\(U=U_{\mu-\mu_0}\) of \cref{lem:ch31-green-kernel} has the following
properties:
\begin{enumerate}
\item\label{itm:ch26-weighted:bfj}
 \(u=U\circ r\) is a continuous BFJ \(\omega\)-psh function and
 \(\MA_\omega(u)=\iota_*\mu\), where \(\iota:\Gamma\hookrightarrow X^{\An}\).
 It is a uniform limit of \emph{real} \(\omega\)-psh model functions
 factoring through \(r\).
\item\label{itm:ch26-weighted:energy}
 \(U\) is absolutely continuous on every edge, with square-integrable
 derivative, and
 \begin{equation}
 \label{eq:ch26-weighted-green-ibp}
 \sum_e c_e\int_e|U'|^2\,\rd s
       =-\int_\Gamma U\,\rd(\mu-\mu_0)<\infty.
 \end{equation}
 For two such measures and potentials, the same identity holds for
 their differences, with \(\mu_1-\mu_2\) on the right.
\end{enumerate}
More generally, if \(U\) is continuous with finite signed distributional
Laplacian, then \(U\circ r\) is BFJ \(\omega\)-psh if and only if
\(\mu_0+\Delta_\Gamma U\geq0\).  In that case its Monge--Amp\`ere
measure is \(\iota_*(\mu_0+\Delta_\Gamma U)\), and the preceding
conclusions hold after adding a normalizing constant.
\end{lemma}

\begin{proof}
\emph{Real model functions on rational subdivisions.}
The original semistable edges have length one in the chosen normalization.
In the convention of \cref{def:ch08-semistable-curve-model}, the components
are smooth and the reduced special fiber is SNC; residue degrees need not
be one.  If the broader semistable convention permits an ordinary
self-node whose two branches are defined over its node residue field,
first blow up each such node.  The exceptional
component has multiplicity two and inserts the midpoint of the
corresponding loop, replacing it by two edges of length \(1/2\).
The strict transform is smooth there.  This preliminary modification
does not change the underlying metric graph or its retraction, and the
nonreduced SNC model is included in the calculation below.  No smooth
point blowups, which can attach new branches, are required here.
No loop description for a conjugate-branch nonsplit node is used.
Any finite set of rational points on them can be inserted by successive
blowups of nodes.  In the local chart \(xy=\pi\), node blowups insert
the sum of the two adjacent primitive valuation vectors.  Repeated sums
give every primitive pair of positive integers by the Euclidean algorithm,
hence every rational point of the interval.  Blowing up the full closed
node performs this construction over its residue field and therefore over
\(K\); a rational residue point is not required.  Finitely many prescribed
points require finitely many blowups.  These are regular SNC models and
may have nonreduced special fiber.  Their underlying skeleton is the same
\(\Gamma\), with a subdivision and the same retraction.

Write the special fiber on such a model \(\mathcal Y\) as
\(\sum_i b_iE_i\).  For arbitrary real values \(a_i\) on its vertices,
the real vertical divisor \(D=\sum_i b_i a_i E_i\) is Cartier with real
coefficients, since \(\mathcal Y\) is regular.  Its model function has
these values and is affine on each subdivided edge; call its skeletal
restriction \(U\).  This follows real-linearly from
\cref{lem:model-function-affine,prop:vertical-divisor-affine-iso}.
Thus every real piecewise affine function with rational breakpoints is an
actual member of \(\mathcal D(X)_{\mathbf R}\).  Neither rational slopes
nor rational coefficients of \(D\) are asserted.

The multiplicities must be retained when computing its curvature.  At
adjacent vertices of multiplicities \(b_i,b_j\), the node subdivisions
above have length \(1/(b_i b_j)\).  Indeed adjacent primitive valuation
vectors have determinant one, preserved when their sum is inserted.
Their component multiplicities are the sums of their coordinates, so
their normalized positions differ by \(1/(b_i b_j)\).
All nodes inserted over an original node retain its residue field.
The node's intersection degree over
\(k\) is \(c_e\).  The algebraic identity
\eqref{eq:ch20-curve-laplacian}, multiplied by \(b_i\), gives
\[
 b_i(D\cdot E_i)
   =\sum_{j\ne i}b_i b_j(E_j\cdot E_i)(a_j-a_i)
   =\sum_{e\ni i}c_e\,\partial_e U(i).
\]
Each node is counted with its residue degree, not merely counted once.
The first equality is valid for real divisors and for the nonreduced
special fibers in question.  The second uses the displayed edge length
and includes a separate summand for each node.
By \eqref{eq:model-mixed-measure}, the curvature coefficient at \(i\)
is exactly \(b_i(D\cdot E_i)\).  The pulled-back reference form has its
old measure \(\mu_0\): new exceptional components have degree zero by
the projection formula.  Consequently
\begin{equation}
\label{eq:ch26-real-weighted-model-measure}
 \MA_\omega(U\circ r)=\iota_*(\mu_0+\Delta_\Gamma U)
\end{equation}
as a signed model measure.  Since each \(b_i>0\), positivity of this
measure is equivalent to nonnegative degree of the real model class on
every special-fiber component, which is relative nefness for a curve
model.  This proves the real model assertion without a line-bundle
positivity promotion or a scalar-pushforward sign argument.

\emph{Arbitrary positive graph measures.}
Choose positive finite atomic measures \(\mu_j\) of mass \(V\), supported
at rational points of the original edges and original vertices, converging
weakly to \(\mu\).  For example, partition every edge into intervals
of mesh tending to zero, assign each interval's mass to a rational point
inside it, and keep vertex atoms at their vertices.
Fix an original vertex \(q\), and normalize the graph potentials
\(U_j=U_{\mu_j-\mu_0}\) and \(U=U_{\mu-\mu_0}\) to vanish at \(q\).
The elementary \cref{lem:ch31-green-kernel} gives uniform convergence
\(U_j\to U\).  Each \(U_j\) is piecewise affine with rational
breakpoints, since its Laplacian is finitely supported there.  The preceding
construction makes \(U_j\circ r\) an actual real model function, and
\eqref{eq:ch26-real-weighted-model-measure} makes it \(\omega\)-psh.
Uniform limits are BFJ-psh by
\cref{thm:ch21-stability}\,\cref{itm:ch21-stab:unif}.
Uniform continuity of the bounded BFJ measure operator,
\cref{prop:ch24-continuous-uniform}\,\cref{itm:ch24-cu:unif}, now yields
\(\MA_\omega(U\circ r)=\iota_*\mu\).
This proves \cref{itm:ch26-weighted:bfj} in the original field and for
the original real form.  The argument has not used the analytic equation
theorem or the line-bundle proposition of \cref{chap:bfj-examples}.

\emph{The energy limit, not just a uniform limit.}
Edgewise integration by parts for the piecewise affine functions gives
\[
 \sum_e c_e\int_e|U_j'-U_k'|^2\,\rd s
 =-\int_\Gamma(U_j-U_k)\,\rd(\mu_j-\mu_k)
 \leq 2V\lVert U_j-U_k\rVert_\infty.
\]
Thus their derivatives are Cauchy in the finite direct sum of weighted
\(L^2\) edge spaces.  On each edge, pass to the limit in
\(U_j(t)-U_j(a)=\int_a^tU_j'(s)\,\rd s\).  Uniform convergence of the
functions and \(L^2\), hence \(L^1\), convergence of the derivatives
show that \(U\) is absolutely continuous with the limiting derivative.
The energies converge.  Also
\(-\int U_j\,\rd(\mu_j-\mu_0)\to-\int U\,\rd(\mu-\mu_0)\)
by uniform convergence and weak convergence with bounded total variation.
This proves \eqref{eq:ch26-weighted-green-ibp}.  Applying the same argument
to simultaneous atomic approximations of two positive measures gives
strong convergence of both derivative sequences and the asserted
difference identity.  A point skeleton has zero energy and is immediate.

\emph{The converse for a given continuous skeletal function.}
Suppose \(u=U\circ r\) is BFJ-psh.  For any real piecewise affine
\(T\) with rational breakpoints, \(t=T\circ r\) is the real model
function constructed above.  Approximate the continuous psh function
\(u\) uniformly by real psh model functions using decreasing BFJ
regularization and Dini's theorem.  Apply model integration by parts to
each approximant and this fixed \(t\), then use uniform continuity of
the BFJ measure and the finite signed model measure \(\ddc t\).  The
result is
\[
 \int_\Gamma T\,\rd\bigl(r_*\MA_\omega(u)-\mu_0\bigr)
 =\int_{X^{\An}}u\,\ddc t
 =\int_\Gamma U\,\rd\Delta_\Gamma T
 =\int_\Gamma T\,\rd\Delta_\Gamma U.
\]
The last equality is precisely the distributional definition, and these
test functions are uniformly dense in \(C(\Gamma)\), by linear
interpolation on finer rational subdivisions.  Hence
\(\mu_0+\Delta_\Gamma U=r_*\MA_\omega(u)\geq0\).
Conversely, positivity allows the preceding construction with
\(\mu=\mu_0+\Delta_\Gamma U\).  Its potential and \(U\) differ by
a constant by the uniqueness part of \cref{lem:ch31-green-kernel}; this
proves the asserted analytic measure identity and sufficiency as well.
\end{proof}

\begin{proposition}
\label{prop:ch26-graph-energy}
Let \(X\) be a smooth projective curve, let \(\mathcal X\) be a regular
semistable model with skeleton \(\Gamma\) equipped with its modulus lengths,
and let \(\omega\) be the semipositive real closed model form of the
standing hypotheses, determined on \(\mathcal X\).  Put
\(\mu_0=\MA_\omega(0)\), of mass \(V=(\{\omega\})\), and use the
residue weights \(c_e\) above.  Then, for every piecewise affine
\(\omega\)-psh function \(u\) on \(\Gamma\), identified with its pullback,
\begin{equation}
\label{eq:ch26-graph-energy}
 E_\omega(u)=\int_\Gamma u\,\rd\mu_0
             -\frac12\sum_e c_e\int_e|u'|^2\,\rd s .
\end{equation}
More generally, for two such functions \(u,v\),
\begin{equation}
\label{eq:ch26-graph-energy-difference}
 E_\omega(v)-E_\omega(u)
 =\int_\Gamma(v-u)\,\rd\mu_0
  -\frac12\sum_e c_e\int_e\bigl(|v'|^2-|u'|^2\bigr)\,\rd s .
\end{equation}
Both formulas also hold for the continuous graph solutions of
\cref{lem:ch26-weighted-graph-real-bfj}, including nonatomic measures.
\end{proposition}

\begin{proof}
At model level the dimension-one definition \eqref{eq:model-energy} and
\eqref{eq:ch26-real-weighted-model-measure} give
\[
 E_\omega(u)=\tfrac12\left(\int u\,\rd\mu_0+
                          \int u\,\rd\MA_\omega(u)\right)
 =\int_\Gamma u\,\rd\mu_0+\tfrac12\int_\Gamma u\,\rd\Delta_\Gamma u.
\]
Weighted edgewise integration by parts makes the last term
\(-\tfrac12\sum_e c_e\int_e|u'|^2\,\rd s\).
For a continuous graph solution, use the positive real model
approximants in \cref{lem:ch26-weighted-graph-real-bfj}.
The energy converges by \eqref{eq:ch26-energy-lipschitz}, the reference
integral by uniform convergence, and the weighted derivative energy by
the strong \(L^2\) convergence proved in that lemma.  Thus the same
identity holds in the limit.  A piecewise affine psh function belongs
to this class by the last assertion of that lemma.  Subtracting the
finite formulas for \(u\) and \(v\) proves
\eqref{eq:ch26-graph-energy-difference}.  Uniform convergence alone was
not used to infer convergence of derivative energies.
\end{proof}

\begin{example}[A degree-two node detects the missing weight]
\label{ex:ch26-nonsplit-weighted-energy}
Let \(K=\mathbf R((\pi))\), with \(-\log|\pi|=1\), and blow up
\(\mathbf P^1_{\mathbf R[[\pi]]}\) at the closed point of its special
fiber defined by \(x^2+1\).  Its residue field is \(\mathbf C\).
The strict transform \(D\) and the exceptional component
\(E\simeq\mathbf P^1_{\mathbf C}\) have multiplicity one and meet
transversely at this degree-two node.  The model is regular semistable;
its generic fiber is \(\mathbf P^1_K\).  The skeleton has one edge of
modulus length one and weight \(c_e=2\).

Take \(L=\mathcal O(2)\) and the pullback of its standard model as
reference.  Its component degrees, over \(\mathbf R\), are \((2,0)\),
so \(\mu_0=2\delta_{v_D}\).  Since
\(\mathcal O(E)|_D=\mathcal O_D(p)\) has degree two and
\(\mathcal O(E)|_E=\mathcal O_{\mathbf P^1_{\mathbf C}}(-1)\) has
degree minus two over \(\mathbf R\), the component-degree change for
\(\mathcal O(-E)\) is \((-2,2)\).  Its model potential \(u\) has
values \((0,-1)\) at \((v_D,v_E)\): in the chart \(\pi=t z\),
\(E=(t=0)\), a frame of \(\mathcal O(-E)\) is \(t\), and
\(-\log\lVert1\rVert=\log|t|\).
Hence
\[
 \Delta_\Gamma u=-2\delta_{v_D}+2\delta_{v_E},\qquad
 \MA_\omega(u)=2\delta_{v_E}.
\]
The modified reference model has degrees \((0,2)\), so it is nef.
This is therefore a semipositive example, not just a signed-curvature
calculation.  Directly from \eqref{eq:model-energy},
\[
 E_\omega(u)=\tfrac12(0-2)=-1,
 \qquad
 \sum_e c_e\int_e|u'|^2\,\rd s=2.
\]
Omitting the edge weight would give \(-1/2\), which is false.
The choice of \(\mathcal O(-E)\) reverses the signs in the
\(\mathcal O(E)\) normalization example
\cref{ex:ch31-nonsplit-weight-normalization}.
\end{example}
